\chapter{Metodología de desarrollo}

En este capítulo se describe la metodología seguida durante el desarrollo y validación de las distintas implementaciones de FlashAttention. En primer lugar, se presenta el algoritmo de FlashAttention 2 utilizado como referencia. Posteriormente, se detalla la implementación desarrollada directamente en CUDA, incluyendo las principales decisiones de diseño y optimizaciones aplicadas.

A continuación, se describe la implementación realizada con Triton y el proceso de \textit{autotuning} utilizado para seleccionar sus parámetros de ejecución. Finalmente, se presenta la metodología empleada para verificar la correctitud de los resultados y para detectar y depurar errores durante el desarrollo de los kernels.

Los kernels implementados en CUDA y Triton, así como la implementación de referencia basada en SDPA, se evalúan bajo una configuración común: dimensión de cabeza igual a 128, longitud de secuencia múltiplo de 16, un único elemento por batch y una única cabeza de atención. En todos los casos se emplea atención no causal.


\section{FlashAttention 2}

El algoritmo utilizado como punto de partida es FlashAttention 2, propuesto por Tri Dao et al. Su objetivo es calcular la atención exacta evitando materializar completamente la matriz \(QK^T\) en memoria global. Para ello, las matrices \(Q\), \(K\) y \(V\) se dividen en bloques que pueden procesarse utilizando la memoria disponible dentro de la GPU.

El cálculo se realiza recorriendo sucesivamente los bloques de \(K\) y \(V\). Para cada uno de ellos se calcula un fragmento de \(QK^T\), se actualizan de forma incremental los parámetros necesarios para el \textit{softmax} y se acumula directamente el producto correspondiente con \(V\). De esta forma, la matriz de probabilidades de atención nunca necesita almacenarse completamente en memoria global.

El algoritmo de FlashAttention 2 utilizado como referencia puede expresarse mediante el siguiente pseudocódigo.

\begin{algorithm}[H]
\caption{FlashAttention 2}
\label{alg:flashattention2}
\begin{algorithmic}[1]

\Require \(Q,K,V \in \mathbb{R}^{N\times d}\)
\Require \(B_r\): número de filas por bloque de \(Q\)
\Require \(B_c\): número de filas por bloque de \(K\) y \(V\)
\Ensure \(O = \operatorname{softmax}(QK^T/\sqrt{d})V\)

\State \(T_r \gets \lceil N/B_r \rceil\)
\State \(T_c \gets \lceil N/B_c \rceil\)

\State Dividir \(Q\) en bloques \(Q_1,\ldots,Q_{T_r}\), con \(Q_i \in \mathbb{R}^{B_r\times d}\)
\State Dividir \(K\) y \(V\) en bloques \(K_1,\ldots,K_{T_c}\) y \(V_1,\ldots,V_{T_c}\), con \(K_j,V_j \in \mathbb{R}^{B_c\times d}\)

\For{\(i \gets 1,\ldots,T_r\)}

```
\State \(O_i \gets 0 \in \mathbb{R}^{B_r\times d}\)
\State \(m_i \gets -\infty \in \mathbb{R}^{B_r}\)
\State \(l_i \gets 0 \in \mathbb{R}^{B_r}\)

\For{\(j \gets 1,\ldots,T_c\)}

    \State \(S_{ij} \gets Q_iK_j^T/\sqrt{d}\)

    \State \(m_i^{new}
    \gets
    \max\left(
        m_i,
        \operatorname{rowmax}(S_{ij})
    \right)\)

    \State \(P_{ij}
    \gets
    \exp\left(
        S_{ij}-m_i^{new}
    \right)\)

    \State \(\alpha_i
    \gets
    \exp\left(
        m_i-m_i^{new}
    \right)\)

    \State \(l_i^{new}
    \gets
    \alpha_i \odot l_i
    +
    \operatorname{rowsum}(P_{ij})\)

    \State \(O_i
    \gets
    \alpha_i \odot O_i
    +
    P_{ij}V_j\)

    \State \(m_i \gets m_i^{new}\)
    \State \(l_i \gets l_i^{new}\)

\EndFor

\State \(O_i \gets O_i \oslash l_i\)
\State Escribir \(O_i\) en \(O\)
```

\EndFor

\end{algorithmic}
\end{algorithm}

En el pseudocódigo anterior, \(\operatorname{rowmax}\) y \(\operatorname{rowsum}\) calculan respectivamente el máximo y la suma de cada fila. Los vectores \(m_i\) y \(l_i\) almacenan, para cada fila del bloque \(Q_i\), el máximo y el término de normalización acumulados durante el cálculo del \textit{softmax online}. El operador \(\odot\) representa una multiplicación elemento a elemento con difusión sobre las columnas de la matriz correspondiente, mientras que \(\oslash\) representa la división por fila.

Nótese que el algoritmo propuesto por Tri Dao et al. no especifica el nivel de la jerarquía de memoria en el que deben residir los datos.

\subsection{Algoritmo CUDA}

La implementación CUDA distribuye el cálculo de la atención a nivel de \textit{warp}. Cada \textit{warp} es responsable de un bloque de 16 filas consecutivas de la matriz de consultas \(Q\), de forma que un bloque CUDA contiene varios \textit{warps} que procesan de manera independiente distintos bloques de consultas. La posición global del \textit{warp} determina las filas de \(Q\) y de la salida \(O\) que le corresponden, mientras que su identificador dentro del bloque se emplea para cooperar en la transferencia de \(K\) y \(V\) a memoria compartida. De esta manera, cada \textit{warp} mantiene en registros su bloque de consultas y el acumulador correspondiente a 16 filas de la salida.

Las matrices \(K\) y \(V\) se recorren en bloques de \(16\times128\) elementos. Para ocultar parcialmente la latencia de memoria global se utiliza un \textit{pipeline} con múltiples buffers en memoria compartida. Mientras un bloque de \(K\) y \(V\) está siendo consumido por los \textit{warps}, bloques posteriores se transfieren desde memoria global mediante instrucciones asíncronas \texttt{cp.async}. Los buffers se organizan de forma circular, de manera que el buffer utilizado para el bloque \(i\) puede reutilizarse para almacenar el bloque \(i+\texttt{buffer\_num}\) una vez que todos los \textit{warps} han terminado de consumirlo. La carga de cada bloque de \(K\) y \(V\) se realiza de forma cooperativa entre todos los \textit{warps} del bloque CUDA. Como consecuencia, todos ellos avanzan de forma prácticamente \textit{lockstep} a través del bucle principal: antes de consumir un nuevo bloque deben esperar a que las transferencias asíncronas correspondientes hayan finalizado y sincronizarse mediante una barrera a nivel de bloque. De forma análoga, antes de reutilizar un buffer de memoria compartida, todos los \textit{warps} deben haber terminado de consumir su contenido. Esta sincronización garantiza la corrección del acceso cooperativo a memoria compartida, aunque reduce el grado de independencia temporal entre los distintos \textit{warps} del bloque.
La transferencia desde memoria compartida al fichero de registros constituye un segundo nivel de \textit{buffering}. Los fragmentos de \(K\) y \(V\) se cargan mediante \texttt{ldmatrix.sync.aligned.m8n8.x4.shared.b16}, que permite obtener simultáneamente cuatro matrices \(8\times8\) desde memoria compartida y distribuirlas directamente entre los registros del \textit{warp}. El código mantiene un pequeño \textit{pipeline} entre estas cargas y las operaciones matriciales: mientras un par de fragmentos se utiliza en los Tensor Cores, se carga desde memoria compartida el siguiente par que será consumido.

Los productos matriciales se realizan mediante la instrucción \texttt{mma.sync.aligned.m16n8k16} con operandos BF16 y acumulación FP32. Cada instrucción calcula un producto \(16\times16\) por \(16\times8\), produciendo un fragmento \(16\times8\) en precisión simple. Para calcular \(QK^\mathsf{T}\), un bloque lógico de puntuaciones de tamaño \(16\times16\) se representa mediante dos fragmentos \(16\times8\). Tras aplicar el escalado y el \textit{softmax} online, las probabilidades se convierten nuevamente a BF16 y se utilizan como operando izquierdo en los productos \(PV\). El acumulador de salida permanece en FP32 durante todo el recorrido de \(K\) y \(V\), realizándose la conversión a BF16 únicamente al escribir el resultado final.
Para mejorar tanto el acceso a memoria global como el acceso posterior mediante \texttt{ldmatrix}, los bloques de \(K\) y \(V\) se almacenan en memoria compartida utilizando un patrón \textit{swizzled}. La granularidad del \textit{swizzle} es de ocho elementos BF16, equivalentes a 16 bytes. Para una fila \(r\) y un bloque lógico de ocho elementos \(c\), la posición física utilizada en memoria compartida se obtiene mediante

$$
c_{\mathrm{swizzled}} = c \oplus (r \bmod 8).
$$

El XOR redistribuye entre los bancos de memoria compartida los segmentos pertenecientes a filas consecutivas. Al efectuar posteriormente un \texttt{ldmatrix}, cada hilo reconstruye la dirección física aplicando el mismo XOR al bloque lógico que necesita. De este modo, la disposición en memoria global puede mantenerse completamente contigua y coalescente, mientras que la disposición en memoria compartida se adapta al patrón de acceso requerido por \texttt{ldmatrix}.

En conjunto, el kernel emplea dos niveles explícitos de movimiento de datos: un \textit{pipeline} memoria global--memoria compartida mediante \texttt{cp.async}, y un segundo \textit{pipeline} memoria compartida--registros mediante \texttt{ldmatrix}. Sobre los fragmentos residentes en registros se ejecutan las instrucciones \texttt{mma.sync} de los Tensor Cores, manteniendo tanto los estados del \textit{softmax} online como el acumulador de salida en registros durante toda la operación.

\begin{algorithm}[H]
\caption{FlashAttention CUDA}
\begin{algorithmic}[1]

\State Asignar a cada \textit{warp} 16 filas consecutivas de \(Q\)
\State Cargar el bloque de \(Q\) correspondiente en registros
\State Inicializar acumulador \(O \gets 0\), máximo de fila \(m \gets -\infty\) y denominador \(l \gets 0\)

\For{\(b \gets 0\) hasta \(\texttt{buffer\_num}-1\)}
\State Cargar cooperativamente \(K_b\) y \(V_b\) de memoria global a memoria compartida mediante \texttt{cp.async}
\State Aplicar \textit{swizzle} XOR durante la escritura en memoria compartida
\EndFor

\For{cada bloque \(i\) de 16 filas de \(K\) y \(V\)}

```
\State Esperar al buffer correspondiente mediante \texttt{cp.async.wait\_group}
\State Sincronizar todos los \textit{warps} del bloque
\Comment{Los \textit{warps} avanzan prácticamente en \textit{lockstep}}

\State \(S \gets 0\)

\State Cargar primeros fragmentos de \(K_i\) desde memoria compartida a registros con \texttt{ldmatrix}

\For{cada fragmento de dimensión \(16\times16\) en la dimensión de cabeza}
    \State Cargar anticipadamente el siguiente fragmento de \(K_i\) mediante \texttt{ldmatrix}
    \State \(S \gets S + QK_i^\mathsf{T}\) mediante \texttt{mma.sync.m16n8k16}
\EndFor

\State \(S \gets S/\sqrt{d}\)

\State \(m_{\mathrm{new}} \gets \max(m,\operatorname{rowmax}(S))\)
\State \(\alpha \gets \exp(m-m_{\mathrm{new}})\)
\State \(P \gets \exp(S-m_{\mathrm{new}})\)
\State \(l_{\mathrm{new}} \gets \alpha l + \operatorname{rowsum}(P)\)

\State Convertir \(P\) de FP32 a BF16 para utilizarlo como operando de los Tensor Cores

\State \(O \gets \alpha O\)

\State Cargar primeros fragmentos de \(V_i\) desde memoria compartida a registros mediante \texttt{ldmatrix}

\For{cada fragmento \(16\times8\) de \(V_i\)}
    \State Cargar anticipadamente el siguiente fragmento de \(V_i\)
    \State \(O \gets O + PV_i\) mediante \texttt{mma.sync.m16n8k16}
\EndFor

\State \(m \gets m_{\mathrm{new}}\)
\State \(l \gets l_{\mathrm{new}}\)

\State Sincronizar todos los \textit{warps} antes de reutilizar el buffer

\If{existen bloques futuros de \(K\) y \(V\)}
    \State Reutilizar el buffer consumido
    \State Cargar asíncronamente el siguiente bloque \(K_{i+\texttt{buffer\_num}},V_{i+\texttt{buffer\_num}}\)
\EndIf
```

\EndFor

\State \(O \gets O/l\)

\State Redistribuir los elementos de \(O\) entre los hilos del \textit{warp}
\State Escribir \(O\) en memoria global mediante stores coalescentes

\end{algorithmic}
\end{algorithm}

\subsection{FlashAttention: Triton}

La implementación en Triton sigue la misma descomposición algorítmica que las versiones anteriores, aunque expresa el cálculo a un nivel de abstracción superior. Cada instancia del programa Triton procesa un bloque de \(B_r\) filas consecutivas de la matriz de consultas \(Q\) y recorre secuencialmente la dimensión de secuencia de \(K\) y \(V\) en bloques de \(B_c\) filas. De esta forma, cada programa mantiene localmente el acumulador de salida y los parámetros necesarios para calcular el \textit{softmax online}, evitando materializar la matriz completa \(QK^\mathsf{T}\) en memoria global.

Dado que la dimensión de cabeza se fija en \(d=128\), cada bloque de consultas tiene dimensiones \(B_r\times128\), mientras que los bloques de claves y valores tienen dimensiones \(B_c\times128\). Para cada bloque de \(K\) se calcula una matriz temporal de puntuaciones de dimensiones \(B_r\times B_c\). Esta matriz únicamente permanece durante la iteración actual del algoritmo y se utiliza inmediatamente para actualizar el \textit{softmax} y el acumulador de salida.

El producto matricial se expresa mediante la operación \texttt{tl.dot}. Con operandos BF16 sobre arquitecturas que disponen de Tensor Cores, el compilador de Triton puede transformar esta operación en las instrucciones matriciales correspondientes, manteniendo la acumulación en FP32. A diferencia de la implementación CUDA, no se especifican directamente instrucciones como \texttt{cp.async}, \texttt{ldmatrix} o \texttt{mma.sync}, sino que la generación y planificación de estas operaciones se delega al compilador.

La implementación se parametriza mediante cuatro parámetros principales:

\begin{itemize}
\item \(B_r\): número de filas de \(Q\) procesadas por cada instancia del programa.
\item \(B_c\): número de filas de \(K\) y \(V\) procesadas en cada iteración.
\item \texttt{num_warps}: número de \textit{warps} utilizados para ejecutar una instancia del programa.
\item \texttt{num_stages}: número de etapas utilizadas por el \textit{pipeline} de cargas y cómputo generado por Triton.
\end{itemize}

La dimensión de cabeza no forma parte del espacio de búsqueda, ya que se mantiene fija en 128 para todos los experimentos. De forma análoga, el tamaño de batch y el número de cabezas se fijan ambos a uno. Por tanto, el proceso de \textit{autotuning} explora únicamente combinaciones de los cuatro parámetros anteriores.

El valor de \(B_r\) determina la cantidad de filas de salida calculadas simultáneamente y, por tanto, afecta directamente al tamaño del acumulador mantenido por cada programa. El valor de \(B_c\) determina el tamaño de los bloques temporales de puntuaciones y la granularidad con la que se recorren \(K\) y \(V\). Valores elevados de ambos parámetros pueden incrementar la reutilización de datos y el trabajo realizado por programa, pero también aumentan el consumo de registros y memoria compartida. Por su parte, \texttt{num_warps} controla el paralelismo dentro de cada programa, mientras que \texttt{num_stages} modifica la profundidad del \textit{pipeline} utilizado para solapar transferencias de memoria y cómputo.

Las distintas combinaciones se evalúan mediante el mecanismo de \textit{autotuning} de Triton, seleccionándose para cada configuración aquella que presenta el menor tiempo de ejecución. De este modo, la versión Triton permite explorar automáticamente distintas relaciones entre tamaño de bloque, paralelismo y profundidad del \textit{pipeline}, en contraste con la implementación CUDA, donde estas decisiones se controlan explícitamente.

El funcionamiento del kernel puede resumirse mediante el Algoritmo~\ref{alg:flashattention_triton}.

\begin{algorithm}[H]
\caption{FlashAttention Triton}
\label{alg:flashattention_triton}
\begin{algorithmic}[1]

\Require \(Q,K,V\in\mathbb{R}^{N\times128}\)
\Require \(B_r\): filas de \(Q\) procesadas por programa
\Require \(B_c\): filas de \(K\) y \(V\) procesadas por iteración
\Ensure \(O=\operatorname{softmax}(QK^\mathsf{T}/\sqrt{128})V\)

\State Asignar a cada programa un bloque \(Q_i\in\mathbb{R}^{B_r\times128}\)
\State Cargar \(Q_i\)
\State \(O_i \gets 0 \in \mathbb{R}^{B_r\times128}\) en FP32
\State \(m_i \gets -\infty \in \mathbb{R}^{B_r}\)
\State \(l_i \gets 0 \in \mathbb{R}^{B_r}\)

\For{cada bloque \(K_j,V_j\in\mathbb{R}^{B_c\times128}\)}

```
\State Cargar \(K_j\)
\State \(S_{ij} \gets \texttt{tl.dot}(Q_i,K_j^\mathsf{T})/\sqrt{128}\)

\State \(m_i^{new}
    \gets
    \max\left(
        m_i,
        \operatorname{rowmax}(S_{ij})
    \right)\)

\State \(\alpha_i
    \gets
    \exp(m_i-m_i^{new})\)

\State \(P_{ij}
    \gets
    \exp(S_{ij}-m_i^{new})\)

\State \(l_i^{new}
    \gets
    \alpha_i\odot l_i
    +
    \operatorname{rowsum}(P_{ij})\)

\State \(O_i \gets \alpha_i\odot O_i\)

\State Cargar \(V_j\)
\State \(O_i
    \gets
    O_i
    +
    \texttt{tl.dot}(P_{ij},V_j)\)

\State \(m_i \gets m_i^{new}\)
\State \(l_i \gets l_i^{new}\)
```

\EndFor

\State \(O_i \gets O_i\oslash l_i\)
\State Convertir \(O_i\) al formato de salida
\State Escribir \(O_i\) en memoria global

\end{algorithmic}
\end{algorithm}

A diferencia de la versión CUDA, el pseudocódigo no representa de forma explícita el movimiento entre memoria global, memoria compartida y registros. Estas decisiones, junto con la planificación de las operaciones matriciales y el solapamiento entre cargas y cómputo, son realizadas por el compilador de Triton en función de la configuración seleccionada.

\chapter {Entorno de prueba}

